% =========================================================
% 3.13 Counter
% =========================================================

\section{কাউন্টার
(\textenglish{Counter})}

\subsection{কাউন্টার কী?}

যে সিকুয়েন্সিয়াল লজিক সার্কিট ক্লক পালস গণনা করে এবং
নির্দিষ্ট ক্রমে এক অবস্থা থেকে অন্য অবস্থায় পরিবর্তিত হয়,
তাকে \textbf{\textenglish{Counter}} বলে।

সহজভাবে, কোনো ডিজিটাল সার্কিটে কতটি
\textenglish{Clock Pulse} এসেছে তা গণনা করার জন্য
\textenglish{Counter} ব্যবহার করা হয়।

একটি \(3\)-\textenglish{bit Binary Counter} সাধারণত এভাবে গণনা করতে পারে—

\[
000\rightarrow001\rightarrow010\rightarrow011
\rightarrow100\rightarrow101\rightarrow110\rightarrow111
\rightarrow000
\]

অর্থাৎ \(3\)-বিট কাউন্টার মোট \(8\)টি ভিন্ন অবস্থা ধারণ করতে পারে।


% =========================================================
% 3.13.2 Construction and counting capacity
% =========================================================

\subsection{কাউন্টারের গঠন ও গণনা ক্ষমতা}

কাউন্টার তৈরির প্রধান উপাদান হলো
\textenglish{Flip-Flop}।
প্রতিটি \textenglish{Flip-Flop} একটি বাইনারি বিট সংরক্ষণ করে।

সাধারণভাবে \(n\)টি \textenglish{Flip-Flop} ব্যবহার করলে
সর্বোচ্চ—

\[
\boxed{
N=2^n
}
\]

টি ভিন্ন বাইনারি অবস্থা পাওয়া যায়।

এখানে,

\[
n=\text{ফ্লিপ-ফ্লপের সংখ্যা}
\]

এবং,

\[
N=\text{মোট সম্ভাব্য অবস্থা}
\]

কাউন্টারের প্রধান অংশগুলো হলো—

\begin{enumerate}[leftmargin=*,itemsep=2pt]

    \item \textbf{\textenglish{Flip-Flop}:}
    বাইনারি অবস্থা সংরক্ষণ করে।

    \item \textbf{\textenglish{Clock Input}:}
    কখন কাউন্টার পরবর্তী অবস্থায় যাবে তা নিয়ন্ত্রণ করে।

    \item \textbf{\textenglish{Combinational Logic}:}
    প্রয়োজন অনুযায়ী পরবর্তী অবস্থা নির্ধারণ করে।

    \item \textbf{\textenglish{Reset Circuit}:}
    কাউন্টারকে প্রাথমিক অবস্থায় ফিরিয়ে আনে।

    \item \textbf{\textenglish{Output}:}
    কাউন্টারের বর্তমান বাইনারি অবস্থা প্রকাশ করে।

\end{enumerate}

\subsubsection*{উদাহরণ: ৩-বিট কাউন্টার}

একটি \(3\)-বিট কাউন্টারে—

\[
n=3
\]

তাই,

\[
N=2^3=8
\]

অর্থাৎ এটি \(0\) থেকে \(7\) পর্যন্ত মোট \(8\)টি ভিন্ন অবস্থা প্রকাশ করতে পারে।


% =========================================================
% 3-bit State Flow Diagram
% =========================================================

\subsubsection*{৩-বিট কাউন্টারের অবস্থা পরিবর্তন}

\begin{center}
\begin{chThreeDiagram}
\begin{tikzpicture}[
    >=Stealth,
    node distance=1.55cm,
    state/.style={
        draw,
        rounded corners=2pt,
        fill=ICTLightBlue,
        minimum width=1.35cm,
        minimum height=0.75cm,
        align=center
    }
]

\node[state] (s0) {\(000\)};
\node[state, right=of s0] (s1) {\(001\)};
\node[state, right=of s1] (s2) {\(010\)};
\node[state, right=of s2] (s3) {\(011\)};

\node[state, below=1.3cm of s3] (s4) {\(100\)};
\node[state, left=of s4] (s5) {\(101\)};
\node[state, left=of s5] (s6) {\(110\)};
\node[state, left=of s6] (s7) {\(111\)};

\draw[->] (s0) -- (s1);
\draw[->] (s1) -- (s2);
\draw[->] (s2) -- (s3);

\draw[->] (s3) -- (s4);

\draw[->] (s4) -- (s5);
\draw[->] (s5) -- (s6);
\draw[->] (s6) -- (s7);

\draw[->]
(s7.west)
to[out=180,in=180]
node[left,font=\small]{\textenglish{Repeat}}
(s0.west);

\end{tikzpicture}
\end{chThreeDiagram}
\end{center}

\par\begingroup

একটি \(n\)-বিট বাইনারি কাউন্টার সর্বোচ্চ
\(2^n\)টি অবস্থা গণনা করতে পারে।
\par\endgroup



% =========================================================
% 3.13.3 Working
% =========================================================

\subsection{কাউন্টার কীভাবে কাজ করে?}

প্রতিটি সক্রিয় \textenglish{Clock Pulse}-এ কাউন্টার
নির্ধারিত নিয়ম অনুযায়ী তার বর্তমান অবস্থা পরিবর্তন করে।

একটি \(3\)-বিট \textenglish{Up Counter}-এর গণনা সারণি—

\begin{center}
\renewcommand{\arraystretch}{1.2}

\begin{tabular}{|c|c|c|c|c|}
\hline

\rowcolor{ICTLightBlue}

\textbf{ক্লক পালস}
&
\(\boldsymbol{Q_2}\)
&
\(\boldsymbol{Q_1}\)
&
\(\boldsymbol{Q_0}\)
&
\textbf{দশমিক মান}
\\
\hline

0 & 0 & 0 & 0 & 0 \\
\hline

1 & 0 & 0 & 1 & 1 \\
\hline

2 & 0 & 1 & 0 & 2 \\
\hline

3 & 0 & 1 & 1 & 3 \\
\hline

4 & 1 & 0 & 0 & 4 \\
\hline

5 & 1 & 0 & 1 & 5 \\
\hline

6 & 1 & 1 & 0 & 6 \\
\hline

7 & 1 & 1 & 1 & 7 \\
\hline

8 & 0 & 0 & 0 & 0 \\
\hline

\end{tabular}
\end{center}

অর্থাৎ \(111\) অবস্থার পর পরবর্তী ক্লক পালসে আবার—

\[
000
\]

অবস্থায় ফিরে আসে।

তাই এই কাউন্টারটি—

\[
\boxed{
\text{\textenglish{Modulo-8 Counter}}
}
\]

হিসেবে কাজ করে।


% =========================================================
% 3.13.4 Classification
% =========================================================

\subsection{কাউন্টারের শ্রেণিবিভাগ
(\textenglish{Classification of Counter})}

কাউন্টারকে বিভিন্ন বৈশিষ্ট্যের ভিত্তিতে শ্রেণিবিভাগ করা যায়।

\subsubsection*{ক্লক সরবরাহের পদ্ধতি অনুযায়ী}

\begin{enumerate}[leftmargin=*,itemsep=3pt]

    \item \textenglish{Asynchronous Counter}
    \item \textenglish{Synchronous Counter}

\end{enumerate}


% =========================================================
% Asynchronous Counter
% =========================================================

\paragraph{\textenglish{Asynchronous Counter}:}

যে কাউন্টারে সব \textenglish{Flip-Flop} একই বাহ্যিক
\textenglish{Clock} সংকেত সরাসরি গ্রহণ করে না,
তাকে \textenglish{Asynchronous Counter} বলে।

এখানে প্রথম \textenglish{Flip-Flop} বাহ্যিক ক্লক পায় এবং
পরবর্তী \textenglish{Flip-Flop} সাধারণত পূর্ববর্তী
\textenglish{Flip-Flop}-এর আউটপুট দ্বারা ট্রিগার হয়।

একে \textenglish{Ripple Counter}-ও বলা হয়।

\textbf{বৈশিষ্ট্য:}

\begin{itemize}[label={{\bulletfont\textbullet}}, leftmargin=*, nosep]
    \item গঠন তুলনামূলক সহজ।
    \item অতিরিক্ত লজিক কম লাগে।
    \item \textenglish{Propagation Delay} জমা হয়।
    \item তুলনামূলক কম গতির সার্কিটে উপযোগী।
\end{itemize}


% =========================================================
% Synchronous Counter
% =========================================================

\paragraph{\textenglish{Synchronous Counter}:}

যে কাউন্টারে সব \textenglish{Flip-Flop} একই
\textenglish{Clock} সংকেত গ্রহণ করে,
তাকে \textenglish{Synchronous Counter} বলে।

এতে প্রয়োজনীয় \textenglish{Combinational Logic}
প্রতিটি \textenglish{Flip-Flop}-এর পরবর্তী অবস্থা নির্ধারণ করে।

\textbf{বৈশিষ্ট্য:}

\begin{itemize}[label={{\bulletfont\textbullet}}, leftmargin=*, nosep]
    \item সব \textenglish{Flip-Flop} একই ক্লক গ্রহণ করে।
    \item তুলনামূলক দ্রুত কাজ করে।
    \item \textenglish{Ripple Delay} কম।
    \item অতিরিক্ত লজিক সার্কিট প্রয়োজন হতে পারে।
\end{itemize}


% =========================================================
% Conceptual Diagram
% =========================================================

\subsubsection*{\textenglish{Asynchronous} ও
\textenglish{Synchronous Counter}-এর ধারণাগত পার্থক্য}

\begin{center}
\begin{chThreeDiagram}
\begin{tikzpicture}[
    >=Stealth,
    line width=0.8pt,
    ff/.style={
        draw,
        rounded corners=2pt,
        fill=ICTLightBlue,
        minimum width=1.7cm,
        minimum height=1.05cm,
        align=center
    }
]

% -------------------------
% Asynchronous
% -------------------------

\node[font=\bfseries] at (2.8,2.1)
{\textenglish{Asynchronous Counter}};

\node[ff] (a0) at (0,0) {\textenglish{FF0}};
\node[ff] (a1) at (2.8,0) {\textenglish{FF1}};
\node[ff] (a2) at (5.6,0) {\textenglish{FF2}};

\draw[->] (-1.5,0) -- (a0.west);
\node[above] at (-0.9,0.1) {\textenglish{CLK}};

\draw[->] (a0.east) -- (a1.west);
\draw[->] (a1.east) -- (a2.west);

\node[below,font=\small] at (2.8,-0.9)
{\textenglish{Clock ripples from one stage to next}};

% -------------------------
% Synchronous
% -------------------------

\node[font=\bfseries] at (9.5,2.1)
{\textenglish{Synchronous Counter}};

\node[ff, fill=ICTLightGreen] (s0) at (7.0,0) {\textenglish{FF0}};
\node[ff, fill=ICTLightGreen] (s1) at (9.5,0) {\textenglish{FF1}};
\node[ff, fill=ICTLightGreen] (s2) at (12.0,0) {\textenglish{FF2}};

\draw (6.2,-1.2) -- (12.8,-1.2);
\node[left] at (6.2,-1.2) {\textenglish{CLK}};

\foreach \x in {7.0,9.5,12.0}
{
    \draw[->] (\x,-1.2) -- (\x,-0.55);
    \fill (\x,-1.2) circle (1.6pt);
}

\node[below,font=\small] at (9.5,-1.8)
{\textenglish{Common clock to all stages}};

\end{tikzpicture}
\end{chThreeDiagram}
\end{center}


% =========================================================
% By Counting Direction
% =========================================================

\subsubsection*{গণনার দিক অনুযায়ী}

\paragraph{\textenglish{Up Counter}:}

যে কাউন্টার ক্রমান্বয়ে গণনার মান বৃদ্ধি করে,
তাকে \textenglish{Up Counter} বলে।

উদাহরণ—

\[
000\rightarrow001\rightarrow010\rightarrow011
\rightarrow100\rightarrow\cdots
\]

\paragraph{\textenglish{Down Counter}:}

যে কাউন্টার ক্রমান্বয়ে গণনার মান হ্রাস করে,
তাকে \textenglish{Down Counter} বলে।

উদাহরণ—

\[
111\rightarrow110\rightarrow101\rightarrow100
\rightarrow011\rightarrow\cdots
\]

\paragraph{\textenglish{Up-Down Counter}:}

যে কাউন্টার নিয়ন্ত্রণ সংকেত অনুযায়ী ঊর্ধ্বক্রমে বা নিম্নক্রমে
উভয়ভাবেই গণনা করতে পারে,
তাকে \textenglish{Up-Down Counter} বলে।


% =========================================================
% Modulus Based
% =========================================================

\subsubsection*{মডুলাস বা গণনা-ক্রম অনুযায়ী}

\paragraph{\textenglish{Modulo-N Counter}:}

যে কাউন্টার মোট \(N\)টি অবস্থা অতিক্রম করার পর আবার
প্রাথমিক অবস্থায় ফিরে আসে,
তাকে \textenglish{Modulo-N Counter} বলে।

উদাহরণ—

\[
0\rightarrow1\rightarrow2\rightarrow\cdots
\rightarrow(N-1)\rightarrow0
\]


% =========================================================
% Decade
% =========================================================

\paragraph{\textenglish{Decade Counter}:}

\textenglish{Decade Counter} হলো একটি বিশেষ
\textenglish{Modulo-10 Counter}।

এটি—

\[
0\rightarrow1\rightarrow2\rightarrow\cdots
\rightarrow9\rightarrow0
\]

ক্রমে গণনা করে।

এর জন্য সাধারণত \(4\)টি \textenglish{Flip-Flop} প্রয়োজন, কারণ—

\[
2^3=8<10\leq16=2^4
\]

চারটি \textenglish{Flip-Flop}-এ \(16\)টি সম্ভাব্য অবস্থা থাকলেও
\textenglish{Decade Counter} কেবল \(10\)টি অবস্থা ব্যবহার করে।


% =========================================================
% Ring
% =========================================================

\paragraph{\textenglish{Ring Counter}:}

\textenglish{Ring Counter} একটি
\textenglish{Shift Register}-ভিত্তিক কাউন্টার,
যেখানে শেষ \textenglish{Flip-Flop}-এর আউটপুট
প্রথমটির ইনপুটে ফিরিয়ে দেওয়া হয়।

একটি \(4\)-বিট \textenglish{Ring Counter}-এর উদাহরণ—

\[
1000\rightarrow0100\rightarrow0010
\rightarrow0001\rightarrow1000
\]

এখানে একটি \(1\) ধারাবাহিকভাবে এক অবস্থান থেকে অন্য অবস্থানে যায়।


% =========================================================
% Johnson
% =========================================================

\paragraph{\textenglish{Johnson Counter}:}

\textenglish{Johnson Counter} একটি
\textenglish{Shift Register}-ভিত্তিক কাউন্টার,
যেখানে শেষ \textenglish{Flip-Flop}-এর বিপরীত আউটপুট
প্রথম \textenglish{Flip-Flop}-এর ইনপুটে দেওয়া হয়।

একটি \(n\)-বিট \textenglish{Johnson Counter}-এ সাধারণত—

\[
\boxed{
2n
}
\]

টি ভিন্ন অবস্থা পাওয়া যায়।


% =========================================================
% 3.13.5 Comparison
% =========================================================

\subsection{অ্যাসিঙ্ক্রোনাস ও সিঙ্ক্রোনাস কাউন্টারের পার্থক্য}

\begin{center}
\renewcommand{\arraystretch}{1.25}

\begin{tabularx}
{\textwidth}
{|L{0.21\textwidth}|Y|Y|}

\hline

\rowcolor{ICTLightBlue}

\textbf{বিষয়}
&
\textbf{\textenglish{Asynchronous Counter}}
&
\textbf{\textenglish{Synchronous Counter}}
\\
\hline

ক্লক সংযোগ
&
সব \textenglish{Flip-Flop} একই বাহ্যিক ক্লক সরাসরি পায় না।
&
সব \textenglish{Flip-Flop} একই ক্লক সংকেত পায়।
\\
\hline

ট্রিগার
&
এক ধাপের আউটপুট পরবর্তী ধাপকে ট্রিগার করে।
&
সব ধাপ একই ক্লকে কাজ করে।
\\
\hline

গতি
&
তুলনামূলক কম।
&
তুলনামূলক বেশি।
\\
\hline

বিলম্ব
&
\textenglish{Ripple Delay} বেশি।
&
তুলনামূলক কম।
\\
\hline

সার্কিট
&
তুলনামূলক সহজ।
&
তুলনামূলক জটিল।
\\
\hline

ব্যবহার
&
কম গতির সাধারণ গণনা ও
\textenglish{Frequency Divider}-এ।
&
উচ্চগতির ডিজিটাল সিস্টেমে।
\\
\hline

\end{tabularx}
\end{center}


% =========================================================
% 3.13.6 Applications
% =========================================================

\subsection{কাউন্টারের ব্যবহার
(\textenglish{Applications})}

\begin{enumerate}[leftmargin=*,itemsep=2pt]

    \item \textbf{ডিজিটাল ঘড়ি:}
    সেকেন্ড, মিনিট ও ঘণ্টা গণনা করতে।

    \item \textbf{\textenglish{Frequency Counter}:}
    নির্দিষ্ট সময়ে কতটি পালস এসেছে তা গণনা করতে।

    \item \textbf{\textenglish{Frequency Divider}:}
    ক্লক ফ্রিকোয়েন্সি ভাগ বা কমাতে।

    \item \textbf{\textenglish{Event Counting}:}
    কোনো ঘটনা বা বস্তু কতবার ঘটেছে তা গণনা করতে।

    \item \textbf{\textenglish{Timer} ও \textenglish{Countdown Circuit}:}
    নির্দিষ্ট সময় বা গণনা নিয়ন্ত্রণ করতে।

    \item \textbf{ডিজিটাল পরিমাপক যন্ত্র:}
    পালস-ভিত্তিক পরিমাপে।

    \item \textbf{\textenglish{Processor} ও
    \textenglish{Control Circuit}:}
    টাইমিং ও ধাপভিত্তিক কার্যক্রম নিয়ন্ত্রণ করতে।

\end{enumerate}


% =========================================================
% 3.13.7 Advantages and Limitations
% =========================================================

\subsection{কাউন্টারের সুবিধা ও সীমাবদ্ধতা}

\textbf{সুবিধা:}

\begin{itemize}[label={{\bulletfont\textbullet}}, leftmargin=*, nosep]
    \item স্বয়ংক্রিয়ভাবে \textenglish{Clock Pulse} গণনা করতে পারে।

    \item ডিজিটাল গণনা ও টাইমিং ব্যবস্থা সহজ করে।

    \item প্রয়োজন অনুযায়ী বিভিন্ন গণনা-ক্রম তৈরি করা যায়।

    \item \textenglish{Frequency Division} ও
    \textenglish{Event Counting}-এ ব্যবহার করা যায়।

\end{itemize}

\textbf{সীমাবদ্ধতা:}

\begin{itemize}[label={{\bulletfont\textbullet}}, leftmargin=*, nosep]

    \item \textenglish{Flip-Flop} সংখ্যা বাড়লে সার্কিট জটিল হতে পারে।

    \item \textenglish{Asynchronous Counter}-এ
    \textenglish{Ripple Delay} তৈরি হয়।

    \item নির্দিষ্ট \textenglish{Modulo Counter} তৈরিতে
    অতিরিক্ত লজিক প্রয়োজন হতে পারে।

    \item সঠিক \textenglish{Reset} ও প্রাথমিক অবস্থা নির্ধারণ না করলে
    প্রত্যাশিত গণনা-ক্রম পাওয়া নাও যেতে পারে।

\end{itemize}

\par\begingroup

\textbf{মনে রাখবে:}

\[
\boxed{
n\text{টি }\text{\textenglish{Flip-Flop}}
\Rightarrow
2^n\text{টি সর্বোচ্চ অবস্থা}
}
\]

এবং,

\[
\boxed{
\text{\textenglish{Decade Counter}}
=
\text{\textenglish{Modulo-10 Counter}}
}
\]

\par\endgroup

\subsection{কাউন্টারের সার্কিট ও গণনা বিশ্লেষণ}
কাউন্টারের ক্লক সংযোগ ও পরবর্তী অবস্থা নির্ধারণের logic দেখে তার কাজ বোঝা যায়। একই Mod-10 কাজ synchronous বা asynchronous দুইভাবেই করা সম্ভব; শুধু modulus দেখে ক্লক ব্যবস্থার ধরন বলা যায় না। নিচের সার্কিটগুলো আদর্শ flip-flop ও নির্দিষ্ট প্রাথমিক অবস্থার ভিত্তিতে ব্যাখ্যা করা হয়েছে।

\subsubsection*{অ্যাসিঙ্ক্রোনাস Up ও Down counter}
\begin{center}
\begin{minipage}{\linewidth}\centering\begin{chThreeDiagram}\begin{tikzpicture}[>=Stealth,line width=.7pt]
\foreach \i in {0,1,2}{
 \node[chff,minimum height=21mm] (f\i) at (3.3*\i,0) {T FF\\$T=1$\\$\downarrow CLK$};
 \draw[->] (f\i.north)--++(0,.8) node[above] {$Q_\i$};
 \draw (f\i.south)--(3.3*\i,-1.7);\fill (3.3*\i,-1.7) circle(1.4pt);
}
\draw[->] (-2,0)--(f0.west) node[pos=0,left] {CLK};
\foreach \i/\j in {0/1,1/2}{\draw[->] (f\i.east)--(f\j.west) node[midway,above] {$Q_\i$};}
\draw (-1,-1.7)--(6.6,-1.7) node[pos=0,left] {$\overline{CLR}$};
\end{tikzpicture}\end{chThreeDiagram}
\captionof{figure}{নামা প্রান্তে কাজ করা T flip-flop দিয়ে তিন বিটের ripple up counter।}

\end{minipage}
\end{center}
\textbf{Up-এর ব্যাখ্যা:} প্রতিটি T=১, তাই নিজস্ব সক্রিয় clock edge-এ toggle করে। FF0 সরাসরি বাইরের clock পায়। FF1-এর clock হলো $Q_0$, FF2-এর clock হলো $Q_1$। $Q_0$ যখন ১ থেকে ০ হয়, FF1 toggle করে; $Q_1$ যখন ১ থেকে ০ হয়, FF2 toggle করে। ফলে $Q_2Q_1Q_0$-তে $000,001,010,\ldots,111,000$ ক্রম পাওয়া যায়।

\textbf{Down-এর সংযোগ:} একই falling-edge flip-flop-এ পরের clock-এ Q-এর বদলে $\overline Q$ দাও। তখন $000\to111\to110\to101\to\cdots\to000$ হবে। rising-edge flip-flop নিলে সংযোগের নিয়মও সেই edge অনুযায়ী বদলাতে হবে। Q না $\overline Q$—এটি না বলে যেকোনো ripple সংযোগকে Up বলা ভুল।

\begin{center}
\begin{minipage}{\linewidth}\centering\begin{chThreeDiagram}\begin{tikzpicture}[>=Stealth,line width=.7pt]
\foreach \i in {0,1,2}{\node[chff,minimum height=21mm] (f\i) at (3.3*\i,0) {T FF\\$T=1$\\$\downarrow CLK$};\draw[->] (f\i.north)--++(0,.8) node[above] {$Q_\i$};}
\draw[->] (-2,0)--(f0.west) node[pos=0,left] {CLK};
\foreach \i/\j in {0/1,1/2}{\draw[->] (f\i.east)--(f\j.west) node[midway,above] {$\overline Q_\i$};}
\end{tikzpicture}\end{chThreeDiagram}
\captionof{figure}{নামা প্রান্তের flip-flop-এ পরের clock হিসেবে পূরক Q দিয়ে ripple down counter।}

\end{minipage}
\end{center}

\chThreeExample{১}{রিপল বিলম্ব}
প্রতিটি flip-flop-এর propagation delay ৮ ns। তিন বিটের ripple-এ তিন ধাপ বদলাতে সর্বোচ্চ আনুমানিক $3\times8=24$ ns লাগতে পারে।
\chThreeAnswer 011 থেকে 100-তে বদলানোর মধ্যে অস্থায়ী 010 ও 000 দেখা যেতে পারে। স্থির ফল পড়ার আগে settling time দিতে হয়। এই সংখ্যাটি উদাহরণের নির্ধারিত delay; সব IC-এর স্থায়ী মান নয়।

\subsubsection*{সিঙ্ক্রোনাস তিন বিটের Up counter}
T flip-flop-এ $T_0=1$, $T_1=Q_0$, $T_2=Q_1Q_0$ দাও। সব clock এক সাধারণ লাইনে। JK ব্যবহার করলে প্রতিটি ধাপে $J_i=K_i=T_i$ দিলেও একই কাজ হয়।
\begin{center}
\begin{minipage}{\linewidth}\centering\begin{chThreeDiagram}\begin{tikzpicture}[>=Stealth,line width=.7pt]
\foreach \i in {0,1,2}{
 \node[chff] (f\i) at (3.5*\i,0) {T FF\\$Q_\i$};
 \draw (f\i.south)--(3.5*\i,-1.7);\fill (3.5*\i,-1.7) circle(1.4pt);
 \draw[->] (f\i.east)--++(.7,0) node[right] {$Q_\i$};
}
\draw[->] (0,2.5)--(f0.north) node[pos=0,above] {$T_0=1$};
\draw[->] (1,0)--(1,3.3)--(3.5,3.3)--(f1.north);
\node[above] at (2.6,3.3) {$T_1=Q_0$};
\fill (1,0) circle(1.4pt);
\node[and gate US,chgate] (a) at (5.5,2.2) {};
\draw (1,2.5)--(4,2.5)|-(a.input 1);\fill (1,2.5) circle(1.4pt);
\draw (4.5,0)--(4.5,1.6)|-(a.input 2);\fill (4.5,0) circle(1.4pt);
\draw[->] (a.output)--(7,2.2)--(f2.north) node[midway,right] {$T_2$};
\draw (-1,-1.7)--(7,-1.7) node[pos=0,left] {CLK};
\end{tikzpicture}\end{chThreeDiagram}
\captionof{figure}{একই ক্লকে তিন flip-flop; নিচের বিটগুলো ১ হলে উপরের বিট toggle হয়।}

\end{minipage}
\end{center}
\textbf{ব্যাখ্যা:} প্রতি ক্লকে $Q_0$ বদলায়। $Q_0=1$ থাকলে carry হওয়ার কারণে পরের ক্লকে $Q_1$ বদলায়। $Q_1Q_0=11$ হলে $Q_2$ বদলায়। সিদ্ধান্ত নেওয়া হয় clock-এর \emph{আগের} মান দিয়ে; একটি ধাপের output-কে পরের ধাপের clock করা হয়নি।

\chThreeExample{২}{সিঙ্ক্রোনাস 011 থেকে পরের অবস্থা}
\chThreeAnswer $T_0=1,T_1=1,T_2=1$। পরের ক্লকে তিন বিট toggle হয়ে 100। সব flip-flop একই clock পায়; বাস্তব ছোট propagation delay থাকলেও ripple clock chain নেই।

\subsubsection*{সিঙ্ক্রোনাস Down এবং Up-Down counter}
Down-এ নিচের সব বিট ০ থাকলে ওপরের বিট বদলাবে। তাই $T_0=1$, $T_1=\overline Q_0$, $T_2=\overline Q_1\overline Q_0$। আগের synchronous চিত্রে FF1-এর T-তে $\overline Q_0$ এবং AND-এর দুই ইনপুটে $\overline Q_1,\overline Q_0$ দিলে Down circuit হয়।

\begin{center}\begin{minipage}{\linewidth}\centering\begin{chThreeDiagram}\begin{tikzpicture}[>=Stealth,line width=.7pt]
\foreach \i in {0,1,2}{\node[chff] (f\i) at (3.5*\i,0) {T FF\\$Q_\i$};\draw (f\i.south)--(3.5*\i,-1.6);\fill (3.5*\i,-1.6) circle(1.4pt);\draw[->] (f\i.east)--++(.6,0) node[right] {$Q_\i$};}
\draw[->] (0,2.6)--(f0.north) node[pos=0,above] {$T_0=1$};
\node[not gate US,draw,minimum size=9mm] (n) at (2.2,2.6) {};
\node[nor gate US,chgate] (r) at (5.6,2.6) {};
\draw (n.input)--++(-.65,0) node[left] {$Q_0$};
\draw[->] (n.output)--(3.5,2.6)--(f1.north) node[midway,right] {$T_1$};
\draw (r.input 1)--++(-.65,0) node[left] {$Q_0$};\draw (r.input 2)--++(-.65,0) node[left] {$Q_1$};
\draw[->] (r.output)--(7,2.6)--(f2.north) node[midway,right] {$T_2$};
\draw (-1,-1.6)--(7,-1.6) node[pos=0,left] {CLK};
\end{tikzpicture}\end{chThreeDiagram}
\captionof{figure}{NOT ও NOR দিয়ে নিচের বিট শূন্য থাকার শর্ত তৈরি করে synchronous down counter।}
\end{minipage}\end{center}
চিত্রে $Q_0,Q_1$ নামযুক্ত gate input সংশ্লিষ্ট flip-flop-এর Q আউটপুটের সঙ্গে যুক্ত। NOR-এ $T_2=\overline{Q_0+Q_1}=\overline Q_0\overline Q_1$; একই ক্লকে প্রয়োজনীয় বিটগুলো একসঙ্গে toggle হয়।

Up-Down-এ U=১-এ Up এবং U=০-এ Down। তার পূর্ণ next-state logic:
\[\begin{aligned}
T_0&=1,\\
T_1&=UQ_0+\overline U\,\overline Q_0,\\
T_2&=UQ_1Q_0+\overline U\,\overline Q_1\overline Q_0.
\end{aligned}\]
প্রতিটি T-এর দুই AND পদ OR করে flip-flop-এ দাও; সাধারণ clock অপরিবর্তিত।
\begin{center}
\begin{minipage}{\linewidth}\centering\begin{chThreeDiagram}\begin{tikzpicture}[>=Stealth,line width=.7pt]
\foreach \i in {0,1,2}{\node[chff] (f\i) at (4*\i,0) {T FF\\$Q_\i$};\draw (f\i.south)--(4*\i,-1.6);\fill (4*\i,-1.6) circle(1.4pt);\draw[->] (f\i.east)--++(.6,0) node[right] {$Q_\i$};}
\draw[->] (0,2.3)--(f0.north) node[pos=0,above] {$T_0=1$};
\node[xnor gate US,chgate] (x0) at (1.8,3.5) {};
\node[xnor gate US,chgate] (x1) at (4.8,2.5) {};
\node[and gate US,chgate] (a) at (7,3.2) {};
\draw (x0.input 1)--++(-.7,0) node[left] {$U$};\draw (x0.input 2)--++(-.7,0) node[left] {$Q_0$};
\draw (x1.input 1)--++(-.7,0) node[left] {$U$};\draw (x1.input 2)--++(-.7,0) node[left] {$Q_1$};
\draw (x0.output)--(3.1,3.5)--(6,3.5)|-(a.input 1);
\draw[->] (3.1,3.5)--(3.1,1.3)--(4,1.3)--(f1.north) node[midway,right] {$T_1$};\fill (3.1,3.5) circle(1.4pt);
\draw (x1.output)--(6,2.5)|-(a.input 2);
\draw[->] (a.output)--(8,3.2)--(f2.north) node[midway,right] {$T_2$};
\draw (-1,-1.6)--(8,-1.6) node[pos=0,left] {CLK};
\end{tikzpicture}\end{chThreeDiagram}
\captionof{figure}{দুই XNOR ও একটি AND-এ U অনুযায়ী একই counter-এর Up/Down নির্বাচন।}

\end{minipage}
\end{center}

এখানে $T_1=\overline{U\oplus Q_0}$ ও $T_2=\overline{U\oplus Q_0}\,\overline{U\oplus Q_1}$। U=১-এ দুই XNOR সরাসরি Q-এর মান দেয়; U=০-এ তাদের পূরক দেয়। ফলে উপরের Up/Down সমীকরণই পাওয়া যায়। Q-নামযুক্ত input সংশ্লিষ্ট Q output-এর সঙ্গে যুক্ত।

\chThreeExample{৩}{দিক পরিবর্তন}
বর্তমান Q=101। U=১-এ পরের state 110; U=০-এ পরের state 100।
\chThreeAnswer একই stored value-এর সঙ্গে direction control পরের state ঠিক করছে। U-কে clock edge-এর কাছে অস্থিরভাবে বদলালে বাস্তব সার্কিটে timing সমস্যা হতে পারে।

\subsubsection*{Mod-N ও Decade counter-এর reset circuit}
একটি binary counter-এ প্রয়োজনীয় Nটি state-এর পরে reset যোগ করলে Mod-N পাওয়া যায়। binary state encoding-এ ন্যূনতম flip-flop সংখ্যা এমন ক্ষুদ্রতম n, যাতে $2^n\geq N$। তাই $n=\lceil\log_2N\rceil$।

\begin{center}
\begin{minipage}{\linewidth}\centering\begin{chThreeDiagram}\begin{tikzpicture}[>=Stealth,line width=.7pt]
\node[draw,minimum width=42mm,minimum height=22mm,align=center] (c) at (0,0) {চার বিট counter\\$Q_3Q_2Q_1Q_0$};
\node[nand gate US,chgate] (n) at (5,0) {};
\draw[->] (-3.5,0)--(c.west) node[pos=0,left] {CLK};
\draw (c.east)--++(.8,0)|-(n.input 1) node[pos=.65,above] {$Q_3$};
\draw (c.east)++(0,-.65)--++(1.1,0)|-(n.input 2) node[pos=.55,below] {$Q_1$};
\draw[->] (n.output)--++(1,0)--++(0,-2)--(0,-2)--(c.south) node[pos=.65,below] {$\overline{CLR}$};
\end{tikzpicture}\end{chThreeDiagram}
\captionof{figure}{শূন্য থেকে শুরু করা চার বিট counter-এ 1010 শনাক্ত করে active-low clear।}

\end{minipage}
\end{center}
\textbf{ব্যাখ্যা:} এই asynchronous-clear উদাহরণে NAND-এর আউটপুট $\overline{CLR}=\overline{Q_3Q_1}$। ০ থেকে বাড়ার পথে প্রথম $Q_3=Q_1=1$ হয় 1010-তে, অর্থাৎ দশমিকে ১০। clear সক্রিয় হয়ে 0000-তে ফিরে আসে; স্থির গণনা ০–৯, modulus ১০। 1010 অল্প সময়ের reset-detection অবস্থা, অতিরিক্ত স্থির গণনা নয়। অন্য initial state বা অনিয়মিত sequence-এ শুধু দুই বিট শনাক্তের যুক্তি যথেষ্ট নাও হতে পারে।

সিঙ্ক্রোনাস Mod-10-এ 1001-এর পরবর্তী clock-এ 0000 লোড করার logic দেওয়া যায়। এতে reset হওয়ার মুহূর্ত clock দ্বারা নির্ধারিত হয়। ripple reset-এর ক্ষণস্থায়ী glitch এবং বাস্তব clear pulse-এর minimum duration বিবেচনা করতে হয়; এ চিত্র আদর্শ কার্যনীতি দেখায়।

\chThreeExample{৪}{Mod-12-এর flip-flop সংখ্যা}
\chThreeAnswer $2^3=8<12\leq16=2^4$; চারটি লাগে। full state detect করলে $R=Q_3Q_2\overline Q_1\overline Q_0$ দিয়ে 1100 শনাক্ত করে clear করা যায়। ১২টি প্রয়োজনীয় state এবং ১৬টি সম্ভাব্য bit pattern এক নয়।

\subsubsection*{Ring counter-এর circuit}
\begin{center}
\begin{minipage}{\linewidth}\centering\begin{chThreeDiagram}\begin{tikzpicture}[>=Stealth,line width=.7pt]
\foreach \i/\b in {0/3,1/2,2/1,3/0}{
\node[chff] (f\i) at (2.7*\i,0) {D FF\\$Q_\b$};
\draw (f\i.south)--(2.7*\i,-1.6);\fill (2.7*\i,-1.6) circle(1.4pt);}
\foreach \i/\j in {0/1,1/2,2/3}{\draw[->] (f\i.east)--(f\j.west);}
\draw[->] (f3.east)--++(.8,0)--++(0,2)--(-1.6,2)|-(f0.west);
\node[above] at (3.5,2) {$Q_0$ ফিরে যায় $D_3$-তে};
\draw (-1,-1.6)--(8.1,-1.6) node[pos=0,left] {CLK};
\end{tikzpicture}\end{chThreeDiagram}
\captionof{figure}{শেষ Q-কে প্রথম D-তে ফেরত দিয়ে চার বিটের ring counter।}

\end{minipage}
\end{center}
\textbf{ব্যাখ্যা:} parallel initialization-এ 1000 বসাও। চার flip-flop একই clock-এ শিফট করে, শেষের বিট প্রথমে ফিরে আসে। $1000\to0100\to0010\to0001\to1000$। একটি ১ ঘুরছে বলে চারটি ব্যবহারযোগ্য state। 0000 দিয়ে শুরু করলে ১ সৃষ্টি হবে না; counter 0000-তেই আটকে থাকবে। one-hot n-bit ring-এর modulus n, $2^n$ নয়।

\subsubsection*{Johnson counter-এর circuit}
\begin{center}
\begin{minipage}{\linewidth}\centering\begin{chThreeDiagram}\begin{tikzpicture}[>=Stealth,line width=.7pt]
\foreach \i/\b in {0/3,1/2,2/1,3/0}{
\node[chff] (f\i) at (2.7*\i,0) {D FF\\$Q_\b$};
\draw (f\i.south)--(2.7*\i,-1.6);\fill (2.7*\i,-1.6) circle(1.4pt);}
\foreach \i/\j in {0/1,1/2,2/3}{\draw[->] (f\i.east)--(f\j.west);}
\node[not gate US,draw,minimum size=9mm,rotate=180] (n) at (6.4,2) {};
\draw (f3.east)--++(.9,0)--++(0,2)--(n.input);
\draw[->] (n.output)--(-1.6,2)|-(f0.west);
\node[above] at (2.8,2) {$\overline Q_0$ ফিরে যায় $D_3$-তে};
\draw (-1,-1.6)--(8.1,-1.6) node[pos=0,left] {CLK};
\end{tikzpicture}\end{chThreeDiagram}
\captionof{figure}{Johnson counter-এ শেষ বিটের পূরক প্রথম D-তে ফেরত আসে।}

\end{minipage}
\end{center}
\textbf{ব্যাখ্যা:} শেষ Q NOT গেটে উল্টে প্রথম D-তে ফিরে যায়। শূন্য থেকে শুরু করলে $0000\to1000\to1100\to1110\to1111\to0111\to0011\to0001\to0000$। চার flip-flop-এ আটটি state, অর্থাৎ সঠিক initialization-এ modulus $2n$। অন্য bit pattern দিয়ে শুরু করলে অন্য cycle-এ আটকে যেতে পারে।

\chThreeExample{৫}{একই চার flip-flop, ভিন্ন গণনা}
\chThreeAnswer পূর্ণ binary counter-এ ১৬ state; one-hot ring-এ ৪; valid Johnson cycle-এ ৮; decade-এ ১০। শুধু flip-flop সংখ্যা দিয়ে ব্যবহৃত modulus জানা যায় না; feedback/reset ও initial state-ও জানতে হয়।

\chThreeExample{৬}{অনেক পালসের পরে state}
শূন্য থেকে শুরু করা Mod-8 up counter-এ ১৯ পালস এসেছে।
\chThreeAnswer $19=2\times8+3$; দুই পূর্ণ cycle শেষে state $3=011$। তিন বিটে ১৯-এর পুরো সংখ্যা সংরক্ষিত হয়নি; শুধু cycle-এর বর্তমান অবস্থা আছে।

\chThreeExample{৭}{ফ্রিকোয়েন্সি ভাগ}
একটি free-running চার বিট binary ripple counter-এ clock ১৬ kHz।
\chThreeAnswer $Q_0$-এ ৮ kHz, $Q_1$-এ ৪ kHz, $Q_2$-এ ২ kHz এবং $Q_3$-এ ১ kHz। সাধারণ ভাগ $f_{Q_i}=f_{CLK}/2^{i+1}$। এটি full binary cycle ও toggle সংযোগের জন্য; যেকোনো Mod-N counter-এর প্রতিটি output-এ একই নিয়ম সরাসরি বসানো যাবে না।

\chThreeExample{৮}{সেকেন্ড গণনা ও সময়ের শর্ত}
এক সেকেন্ড পরপর একটি নির্ভরযোগ্য পালস এবং Mod-60 counter দিয়ে ০–৫৯ সেকেন্ড গণনা করা যায়।
\chThreeAnswer ছয় বিট লাগে, কারণ $32<60\leq64$। প্রতি ৬০ পালসে পরের minute counter-এ একটি carry/event signal দিতে হয়। প্রতি clock-কে এক সেকেন্ড ধরা যাবে কেবল clock frequency ১ Hz হলে।

% Mahbub teaching-type follow-up, 2026-10-14
\chThreeExample{৯}{counter-এর output থেকে pulse সংখ্যা}
শূন্য থেকে শুরু করা তিন-বিট Up counter-এর বর্তমান state 101। মোট কয়টি pulse এসেছে?
\chThreeAnswer state-এর মান ৫। প্রথম cycle-এর মধ্যে হলে pulse সংখ্যা ৫; cycle-এর সংখ্যা জানা না থাকলে $N=8k+5$, যেখানে $k=0,1,2,\ldots$। ফলে ৫, ১৩, ২১ ইত্যাদি pulse-এর পর একই state আসে। বর্তমান state দেখে মোট elapsed pulse সবসময় এককভাবে নির্ণয় করা যায় না।

